\section{Low cone residues for a genuine first native two-cell}
\label{sec:136-first-cell-residues}

\paragraph{Archive and scope.}
This concerns cote136, physical p.~117/footer116/scan117,
physical p.~119/footer118/scan119, and the common-source question
physical p.~121/footer120/scan121. The old algebra in this result
is the actual left Kan envelope \(L(F)\) of the common first
skeleton. The new native input has a prescribed, possibly nonzero,
full boundary. The conclusion includes the inherited RH1 iteration
and, by the continuation proved below, subsequent free horizontal
cell attachments. It does not include further positive-form vertical
repairs. It does
not assume a simplicial retraction of the new algebra onto \(L(F)\).

We use the checked inputs ARCH-136-RELATIVE-DISK-FACTORIZATION,
ARCH-136-HIGH-NATIVE-RESIDUE, ARCH-136-NATIVE-ALL-ARITIES,
ARCH-136-TAIL-EXT-VANISHING, and the full supplied
AQ-136-RELATIVE-H1. The complete native pushout is the one
defined in ARCH-136-RELATIVE-CELL-PATHS.

Let \(S=\Gamma_{\mathbb Z}(T)\). Fix finitely many admissible
targets \(C_j\) with the literal common first skeleton \(F\),
where \(F_0=S\) and \(F_1\) is the specified free algebra
on even zero-form unit disks. Put \(A=L(F)\). Thus each
\(A_r\) is a tensor product of free even zero-form unit disks
over \(S\), and the canonical maps \(A\to C_j\) are the
identity on \(F\). Let \(p\) be prime, \(q\geq1\) odd,
\(m>0\), \(D=p(q+1)\), \(v=pm\), and
\[
                    X=\Delta[2],\qquad B=\partial X
                                      =\operatorname{sk}_1X.
                                                               \tag{136.FC1}
\]
Prescribe a native boundary section \(a_B\in F^{m,q}(B)\).
Prescribe complete coherent boundary parameters on a free
resolution of \(I_M\), \(M\geq1\), with generator weights
at least \(M\) and with both bottom relations, so that
there is a strict map \(\mathcal N(B)\to A\). Form
\[
                   E=A\amalg_{\mathcal N(B)}\mathcal N(X).
                                                               \tag{136.FC2}
\]
Suppose strict maps \(f_j:E\to C_j\) are given extending
the canonical old maps. Their coherent parameter images may
already be prescribed.

The boundary-parameter hypothesis is available here without
adjoining old higher forms: on each edge, contract
\(z_B=\gamma_p(da_B)\) and
\(e_B=a_B\gamma_{p-1}(da_B)\) in the free even-disk
algebra \(F_1\): first take a primitive \(c_B\) of
\(z_B\), then a primitive \(t_B\) of the closed element
\(pc_B-e_B\). Thus there are common \(c_B,t_B\) in forms
\(D-1,D-2\) with \(dc_B=z_B\), \(dt_B=pc_B-e_B\).
These positive forms agree at vertices. Taking
\(U_0=\epsilon c_B\), \(V_0=\epsilon t_B\), and all
higher parameters zero satisfies the full boundary equations.
The common-skeleton coherent-lift theorem supplies target
extensions after the usual input-row and parameter-tail guards.
The result below also permits any other already given complete
boundary and endpoint parameter data.

Choose a bound \(W(D)\) dominating the finitely many target
comparison bounds through form degree \(D+1\), and the
endpoint-filling bound of \(F_1\). Impose the stronger,
original-native-weight condition
\[
                              v\geq W(D).                \tag{136.FC3}
\]
It is not enough here to increase only \(M+v\).

\paragraph{ARCH-136-FIRST-CELL-RESIDUES.}
The actual integral source (136.FC2) admits the supplied RH1
iteration under \(E\) and over all targets. Its colimit \(E^+\)
still has literal first skeleton \(F\). There is one bound
\(W_{\mathrm{low}}\), depending only on the targets' low
vertical and horizontal cutoffs and the endpoint bound of \(F\),
such that the following holds uniformly in the simplex dimension,
the modulus, and the RH1 stage, including its colimit.

If \(w\geq W_{\mathrm{low}}\), \(N\geq1\), and
\[
 b\in(E^+)_r^{w,2},\qquad z\in(E^+)_r^{w,3},
                         db=Nz,\qquad dz=0,              \tag{136.FC4}
\]
then there are target primitive pairs
\[
 dc_j=f_j(z),\qquad dt_j=Nc_j-f_j(b),
                  \qquad |c_j|=2,\quad |t_j|=1,          \tag{136.FC5}
\]
whose restrictions to \(G_r=\operatorname{sk}_1\Delta[r]\)
are identical as sections of the common graph algebra \(F(G_r)\).
Thus every degree-three graph-difference residue of this genuine
first-cell source vanishes, including its nonlinear residues after
RH1. A primitive in the source itself is not asserted.

\begin{proof}
We first work over \(R=\mathbb Z_{(p)}\), using subscripts
\(R\) only when necessary. Fix a simplicial degree \(r\).
The actual pushout has the pointwise description
\[
              E_r=A_r\otimes_S\bigotimes_{\sigma}\mathcal N_\sigma,
                   \qquad \sigma\in X_r\setminus B_r.  \tag{136.FC6}
\]
The tensor products denote the strict PD coproducts with their
induced differentials. All new blocks are independent over
\(A_r\); their nonzero faces in the old algebra remain part
of the simplicial structure. In particular (136.FC6) does not
supply a simplicial augmentation onto \(A\).

In each native block assign arity one to its input disk and
arity \(p\) to its coherent parameters. Its arity-\(p\)
linear complex \(K_\sigma\), including the native
\(e,z\), is the reversed length-\(D\) truncation of a
free resolution of
\(S/(p,I_M)\langle v\rangle\). It is exact below form
\(D\), its positive form modules are free, and every
generator weight is at least \(v\). These are the checked
native-complex identities; in particular they include the
nonfree bottom syzygy, rather than replacing it by a free module.

We claim the maps of this linear complex to all targets have
nullhomotopies with common graph values. Complete its reversed
truncation to a free resolution \(Q_\bullet\) of the displayed
finite quotient, without introducing generator weights below
\(v\). ARCH-136-TAIL-EXT-VANISHING and the multiplication-by-
\(p\) sequence for the integer-free module \(S/I_M\) give
\[
 H^i\operatorname{Hom}_{S,\mathrm{gr}}(Q_\bullet,S)=0
                                      \quad\text{for all }i.
                                                               \tag{136.FC7}
\]
On the graph \(G_r\), the restrictions of the linear maps
are one common map \(g:K_\sigma\to F(G_r)\), because
every \(f_j\) is the identity on the actual first skeleton.
Apply total Hom from \(Q_\bullet\) to the graph-incidence
diagram of \(F_1\) and its two coefficient endpoints. Its
restriction map is degreewise onto: (136.FC3) places every
free-coordinate weight above the endpoint bound. The integral
coefficient contraction \(F_1\simeq S\) compares its
homotopy fiber with ordinary graph cochains with coefficient
complex \(\operatorname{Hom}(Q_\bullet,S)\). By
(136.FC7) that complex is acyclic. This is the same exact
incidence argument as ARCH-136-SKELETON-COHERENT-MOORE;
here the coefficient complex is zero in cohomology, rather
than concentrated in degree zero.

The degree-\(D\) cocycle represented by \(g\) consequently
has a degree-\(D-1\) primitive. In cochain notation this is
a common graph homotopy
\[
 g_q=d h_{G,q}+h_{G,q+1}d_K\ (q>0),\qquad
                            g_0=h_{G,1}d_K.              \tag{136.FC8}
\]
Only free positive modules of \(K_\sigma\) occur in its
components. Extend those components independently to each
target simplex using the acyclic rows in (136.FC3). The defect
of an extension is a relative total degree-\(D\) cocycle
on \((\Delta[r],G_r)\). The generalized derived comparison
proved in ARCH-136-HIGH-NATIVE-RESIDUE, together with (136.FC7),
makes that relative cohomology group zero. Subtract a relative
primitive. This yields homotopies \(h_j\) of the actual
linear maps with \(h_j|_{G_r}=h_G\). As in that theorem,
the equation on the free cover of \(K_\sigma^0\) descends
to the literal bottom equation in (136.FC8). No extension of
a map out of a nonfree syzygy has been assumed.

Use the same universal free unit-disk factor \(D_\sigma\)
for all these homotopies. On a basis element \(x\in K^a\),
\(a>0\), its generators satisfy \(du_a(x)=v_a(x)\), and
the source map is
\[
 j_0(x)=u_1(d_Kx),\qquad
 j_a(x)=v_a(x)+u_{a+1}(d_Kx)\quad(a>0).                   \tag{136.FC9}
\]
Evaluation in target \(j\) sends \(u_a(x)\) to
\(h_{j,a}(x)\) and \(v_a(x)\) to \(d h_{j,a}(x)\).
All these evaluations agree on the graph, by (136.FC8).
Define the pointwise auxiliary factor
\[
                    D_r=A_r\otimes_S\bigotimes_\sigma D_\sigma.
                                                               \tag{136.FC10}
\]
It is itself a free algebra on unit disks, since \(A_r\)
has that property. The explicit disk calculation gives
\[
                   H^0(D_r)=S,\qquad H^1(D_r)=H^2(D_r)
                                            =H^3(D_r)=0. \tag{136.FC11}
\]
Indeed even lower-degree disks contract integrally. For an odd
lower degree the first noncontractible block starts in form three
with an injective differential to form four. Products of two such
blocks start in form six. This proof works in every weight and
for arbitrary ranks. The target evaluations of \(D_r\) agree
on \(G_r\), including their evaluations on its old factor \(A_r\).

We next pass from the linear complexes to the whole native
algebras, without omitting any arity. Use a separate arity coordinate
for each \(\sigma\) in (136.FC6), and give \(A_r\) residue
zero. The component with all arity coordinates divisible by \(p\)
is
\[
                  E_r^{(0)}=A_r\otimes_S
                         \Gamma_{\mathrm{gr}}(\bigoplus_\sigma K_\sigma).
                                                               \tag{136.FC12}
\]
This identity uses the unit coefficients
\((ps)!/(s!(p!)^s)\) over \(\mathbb Z_{(p)}\), as
proved in ARCH-136-NATIVE-ALL-ARITIES. Its target maps factor
through (136.FC10), by (136.FC9).

Every other residue component is actually contractible as a
differential module over (136.FC12). On a nonzero residue
coordinate its old disk contraction is
\[
 h(\gamma_s(b))=a\gamma_{s-1}(b)/s,\qquad
 h(a\gamma_{s-1}(b))=0,\qquad p\nmid s.                  \tag{136.FC13}
\]
It is linear, with the graded sign, over the subalgebra generated
by \(\gamma_p(b)\) and \(a\gamma_{p-1}(b)\); the
binomial identity in the all-arity theorem verifies this for
every divided power. It extends over all coherent parameters,
over the other blocks, and over \(A_r\). Such a residue
component has no form-zero elements, and hence no one-cycles.

For the source before RH1, decompose a closed degree-two Moore
cone pair into these finitely many residue components. On a
nonzero component, the actual source contraction gives a source
cone primitive (the cone homotopy is \((x,y)\mapsto(hx,-hy)\)).
Its images have common graph restrictions because the source
maps agree on \(F\). The zero-residue component factors
through \(D_r\), whose degree-two Moore-cone cohomology is
zero by (136.FC11), for every modulus. A primitive there has
common graph images under its evaluations. Summing these primitive
pairs proves the assertion locally at \(p\) before RH1.

We verify that this argument survives the actual RH1 construction.
All zeroforms, and all one-cycles, lie in the zero-residue part.
At a fixed simplex the RH1 pushout therefore adjoins only zeroforms
to that subalgebra; all its identifying relations are zero-form
relations there. The other residue modules undergo ordinary module
base change over this enlarged subalgebra. More explicitly, the
full PD ideal is generated by the divided powers of those even
zero-form relations. For their multiples by other elements the
scalar PD identity keeps every additional relation in the module
ideal generated from the same zero-residue relations. Thus the
module homotopies (136.FC13) extend by base change, with cancellation
of the terms differentiating the new coefficient algebra. No
flatness is needed. These components still have no zeroforms and
no one-cycles at every stage.

It remains to retain the factorization of the zero-residue part
through the same \(D_r\). This can be done pointwise, and
requires no simplicial structure on the auxiliary factors.
Let \(V\) be the replacement zero-form module in one RH1
step. In every repaired weight its defining normalized modules
give the exact sequence of simplicial modules
\[
                 0\longrightarrow\underline{S^w}
                   \longrightarrow V^w\xrightarrow{d_V}Z^1(E^{w,*})
                   \longrightarrow0.                    \tag{136.FC14}
\]
Here normalized degree zero of \(Z^1E\) is zero, since
\(E_0=S\); in positive normalized degrees \(d_V\) is
the identity onto the old one-cycles. Exactness follows from
Dold--Kan. The potentially nonzero horizontal boundary
\(V_1\to V_0=S\) is retained in (136.FC14); it has not
been set to zero. Evaluation at one chosen vertex splits the
sequence as modules in the fixed simplicial degree \(r\).

The common graph evaluation of \(D_r\) gives a coefficient-
preserving augmentation \(\eta_r:D_r\to S\) by that
vertex. It need not be the augmentation setting all zero-form
disk generators to zero. For \(x\in V_r\), map it to
the unique element \(\widetilde x\in D_r^0\) satisfying
\[
       d\widetilde x=j(d_Vx),\qquad
       \eta_r(\widetilde x)=\operatorname{ev}_{\mathrm{vertex}}(x).
                                                               \tag{136.FC15}
\]
Existence follows from \(H^1(D_r)=0\), and uniqueness from
\(H^0(D_r)=S\). The assignment is linear and preserves
weight. For an old zeroform both conditions hold for its existing
image in \(D_r\), so all identifying relations of the RH1
PD pushout are respected. The factorization therefore extends.
In each target, its value on a new zeroform and the retained RH1
value have the same differential and the same vertex value.
Their difference is a closed zeroform. In the fixed repair range
the target's closed zeroforms are precisely its coefficients,
so that difference is zero. Thus these are factorizations of the
actual retained simplicial target maps, even though the auxiliary
\(D_r\) was chosen only pointwise. Lower weights are unchanged.

The same argument applies at every stage. At a limit it applies
by the compatible maps already constructed, since each expression
has finite support. Together with the extended contractions, it
proves the asserted common-graph primitive property over
\(\mathbb Z_{(p)}\) after the entire RH1 iteration, in fact
in every weight there.

We now justify the iteration and descend the residue conclusion
integrally. The old algebra and each full native coherent block
are rationally free unit-disk algebras. The checked triangular
change in the coherent presentation, including elimination of
the bottom relation, proves this with arbitrary ranks. Thus
\(E_0=S\), weight zero is \(\mathbb Z[0]\), and the
rational vertical coefficient comparison holds in every simplex
and weight. These are exactly the supplied RH1 hypotheses once
its one fixed low target bound is chosen. The repair keeps the
first skeleton literally \(F\), because there the zero-to-one
differential is already the required normalized isomorphism onto
one-cycles. Its rational transitions remain isomorphisms.

For any prime \(\ell\ne p\), the same triangular coordinate
change is already integral over \(\mathbb Z_{(\ell)}\):
it uses only division by \(p\). In that localization every
coherent parameter part is a free unit-disk extension of the old
native unit disk. Consequently each \(E_{r,(\ell)}\) is
a tensor of free unit-disk algebras, and has the cohomology
(136.FC11) in degrees zero through three in every weight.
RH1 is the identity on such an algebra: its map from old zeroforms
to \(V\) is an isomorphism in every repaired normalized degree.
Moreover the integral RH1 pushout commutes with localization.
Normalization is a direct summand, flat localization commutes
with cycles, its coefficient boundary is defined by the same
unique rational primitive, and the free PD universal property
commutes with scalar extension. Hence
\[
                     (E^+)_{(\ell)}\cong E_{(\ell)}
                             \qquad(\ell\ne p).         \tag{136.FC16}
\]
The same assertion holds at every intermediate RH1 stage.
Every closed low cone pair then has an actual source primitive
locally at \(\ell\ne p\), and its target images agree on
the graph. At \(p\) we proved the common-graph primitive
property directly, without requiring a source primitive.

Finally choose \(W_{\mathrm{low}}\) to ensure the targets'
vertical comparison through degree three, their row extension
in forms zero and one, and the graph coefficient comparison.
It does not depend on \(r,N\) or the stage. For (136.FC4)
choose any integral primitive pair in each target, using
\(H^2(C_j)=H^3(C_j)=0\) in that range. The difference of
two graph restrictions is a closed degree-one Moore-cone element
of \(F(G_r)\). Its cohomology class is independent of the
target primitive choices: their differences are closed degree-one
cone elements in a target simplex, whose cone cohomology is zero
because \(H^1(C_j)=H^2(C_j)=0\).

The graph comparison identifies this difference class with an
element of \(H^1(G_r;S^w/NS^w)\). Restriction to the finite
graph and all these finite-degree kernels commute with flat
localization. The class is zero after localization at every prime,
by the two local conclusions just proved. It is therefore zero
integrally. Choose a graph cone primitive of each difference,
extend its form-zero and form-one components to the corresponding
target using the acyclic rows, and subtract its differential from
the original primitive. All graph restrictions now agree.
Taking the signs \(\beta_j=(-t_j,-c_j)\) gives exactly
(136.FC5). This proves the integral statement with its uniform
bound and its asserted RH1 scope.
\end{proof}

\paragraph{Horizontal continuation of ARCH-136-FIRST-CELL-RESIDUES.}
The conclusion remains valid through any sequence of free horizontal
unit-disk attachments of simplex dimension at least two, interspersed
with RH1 steps using the same fixed low-degree bound, and through
their filtered colimit. In particular one can construct an actual
source \(E\to E^{\mathrm h}\), over all the targets and with
first skeleton \(F\), such that
\[
\begin{gathered}
 H^0((E^{\mathrm h})_r^{w,*})=S^w,\qquad
 H^1((E^{\mathrm h})_r^{w,*})=0\quad(w\geq W_{\mathrm{low}}),\\
 H_s\bigl(N(E^{\mathrm h})_*^{w,a}\bigr)=0
                      \quad(w\geq\mu(a),\ s\geq0),
\end{gathered}                                             \tag{136.FC17}
\]
and every cone pair (136.FC4) in this source admits (136.FC5)
with common graph restrictions. The bounds are fixed throughout.
This is still a partial source: no higher vertical vanishing is
asserted by (136.FC17).

\begin{proof}
At a fixed simplex a free horizontal cell attachment is a tensor
extension by free vertical unit disks. The boundary-module inclusion
splits degreewise, and its splitting tensored with the vertical
disk respects the form differential, exactly as in the supplied
PARTIAL-SOURCE proof. Consequently it sends the local factor
\(D_r\) of the preceding proof to its tensor product with
these same independent free disks. Map the new factor generators
to the already chosen images of the actual new source generators
in each target. Their restrictions to \(G_r\) are common:
the actual source first skeleton is still \(F\), and all its
target maps are the identity there. Thus this is again a factor
with common graph evaluations and the cohomology (136.FC11).

Give the newly added disks residue zero relative to the old
native coordinates. Each old nonzero residue module undergoes
base change by these disks, so its module contraction persists,
and it still has no zeroforms or one-cycles. This argument is
pointwise. It does not require the chosen residue-zero subalgebras
to be preserved by all simplicial operators: the next RH1 step
only uses the pointwise facts that all zeroforms and all one-cycles
lie in that part, and that its map to \(D_r\) factors the
actual retained simplicial target maps. Equations (136.FC14)--
(136.FC15) therefore extend the factor through the next RH1 step
without a new choice of target values. The contractions extend
by the same zero-form PD base-change argument. Induction proves
the common-graph primitive property at the prime \(p\).

At every other prime the localized source starts as a tensor
of free unit disks. A horizontal step adds such disks, and an
RH1 step remains the identity there. Thus its cohomology through
degree three remains (136.FC11). Integral vanishing of the
canonical graph-difference class follows from the same localization
argument as before. At a limit the auxiliary factors and their
maps are compatible unions of the newly adjoined free disks;
the residue-module contractions commute with base change. Elements
have finite support, so both the factor argument and the exact
filtered-colimit cohomology argument remain valid. No pointwise
factor is being claimed to form a simplicial algebra of its own.

For the final assertion, fix the target horizontal cutoffs
\(\mu(a)\) once. In form zero enlarge \(\mu(0)\) above
the bound where the common edge already has boundary \(T^{[w]}\).
Then horizontal degree-zero homology is already zero in the chosen
tails: the coefficient row is filled by \(F_1\), and positive
form rows have simplex-zero term zero. All remaining horizontal
cycles can therefore be killed by cells of dimension at least two.
Alternate the step attaching all such free cells with RH1, always
using the same cutoffs. Rational vertical coefficient exactness
persists, so every RH1 application is legitimate. Every old
horizontal class and every old high-weight vertical one-class
is killed at a later step; zero-classes become coefficients.
The usual finite-support and filtered-colimit argument proves
(136.FC17). The proved residue property holds throughout that
same construction.
\end{proof}

\paragraph{Remaining limitation.}
This handles a genuine boundary-constrained first two-cell,
including nonzero old graph data, arbitrary prescribed coherent
target images, all native arities, and the entire RH1 repair.
The horizontal continuation retains the pointwise unit-disk factor
and its contracted complementary modules. Arbitrary positive-form
vertical repairs have not been included in this theorem; they may
require compatible new primitive choices that these pointwise
factors do not by themselves provide. For all
native primes on one fixed input, the conditions
\(pm\geq W_C(p(q+1))\) also do not follow from a single
finite bound on \(m\). The full integral common-source question
remains unresolved.

\paragraph{Verification status.}
The author completed a whole saved-proof read, including the
actual boundary pushout, common graph nullhomotopy from the full
finite-quotient resolution, residue-component contractions,
the nontrivial horizontal boundary in the RH1 replacement module,
vertex-normalized factor extension, localization of the actual
RH1 pushout, and descent of the canonical graph-difference class.
The added continuation has also been read in full, including the
non-simplicial pointwise residue bookkeeping, actual graph evaluations
of the new horizontal generators, unchanged fixed cutoffs, and
the complete filtered-colimit argument.
Independent saved-proof reviews are recorded separately in the
accompanying event register.
